\section{Introduction}\label{sec:short_intro}

What matters more when pricing an option on the average price of crude
oil: allowing for uncertainty about volatility, or choosing a
volatility estimate informed by option prices? WTI Average Price
Options (APOs) hedge exposure to an arithmetic average of crude-oil
futures settlements, rather than a single price at expiry.
Their value depends on both the pricing model and its volatility
input. Bayesian methods account for uncertainty about that input,
but greater computational complexity does not necessarily produce
a more accurate conditional price.

Asian options depend on a price average rather than a single terminal
price. Geometric-average approaches and conditioning approximations
are long-established tools for computing arithmetic-average
option values \citep{kemnavorst1990,curran1994}. REMEF has
published related Asian-option work: \citet{matiasgonzalez2019}
study a stochastic optimal consumption/investment framework with
Asian-option valuation, whereas \citet{ortizjimenez2026} apply
realized volatility and mean reversion to the CEMEXCPO setting.
Neither compares the effect of averaging over a volatility posterior with the effect of using different volatility inputs to price WTI average-price options.

Bayesian option pricing is also well established. \citet{guidolin2003}
study learning and option-price dynamics, while
\citet{romboutsstentoft2014} report a small effect of parameter
uncertainty on pricing errors in their equity-option setting.
More recently, \citet{sagorbooneghanam2026} propagate Bayesian
uncertainty into illustrative European-option prices, including
a WTI example. Those studies do not make the same comparison for WTI
average-price options. Our contribution is not a new method of
posterior integration. It is a direct comparison of posterior
integration, earlier option-implied volatility, and a simple
benchmark that carries forward the last observed implied volatility.

WTI APOs provide a different test setting. Market-specific
models may account for stochastic convenience yields, jumps
or richer futures dynamics \citep{schwartz1997,shirayatakahashi2011,ewaldwu2023}.
We hold the one-factor option-pricing benchmark fixed in
the original analysis, so the influence of volatility
information is not conflated with changes in the dynamics
of the commodity futures curve. This comparison focuses on the volatility input rather than
on selecting a more complex model for crude-oil futures.

The distinction between historical and option-implied information
is important. \citet{gildertsiaras2020} study forecasts of
\emph{subsequent realized crude-oil volatility} extracted from
option prices. We study errors in pricing WTI
\emph{average-price options} at later valuation dates. Volatility inferred from option prices need not match
volatility estimated from historical futures returns,
particularly when the two estimates describe different
underlying quantities. Consequently, a large historical-versus-implied
pricing difference by itself neither identifies a variance
risk premium nor demonstrates that Bayesian inference fails.

We separate three questions. First, how much does averaging
option prices over the volatility posterior change the price itself?
We call this posterior integration (PI), and compare it with pricing
at the posterior mean (PM). Second, how much does the chosen
volatility input affect pricing errors? Third, how accurately can
we describe errors relative to observed settlement prices?
The PI--PM difference, the spread of model prices across
the posterior, and an interval for settlement-price errors
are different quantities; none should stand in for another.

We study these questions using 10 million simulated return
histories and two separate analyses of market data. Each
analysis has its own pricing model and information cutoff. The first uses individual Barchart
APO histories, a CL futures-curve-based one-factor
Asian option approximation, 310 October-expiry
contract-date observations in August--September 2026,
2,381 date-ordered conditional repricing observations across 15 dates,
and a 253-observation cross-instrument check using
earlier LO vanilla-WTI implied volatilities. The second
uses Databento AAO settlement records and a \emph{different},
simplified pricing model that treats the terminal arithmetic
average as lognormal. It compares prices with 97,341
observations classified by the provider as actual settlements
across 189 dates. It also examines Bayesian estimates based on
earlier implied volatility and intervals for subsequent
settlement-price errors.

The results point to a practical distinction. In the original
study, integrating over historical volatility uncertainty changes
the point price very little. Using earlier option-implied
volatility instead of the tested historical estimates yields
much smaller errors relative to settlement-price proxies.
The separate, longer AAO study finds a similar ranking of
volatility inputs, but Bayesian updating of earlier implied
volatility does not reliably outperform carrying forward
the most recent observation. Its pricing-error intervals also fall below their stated
coverage levels in parts of 2026. These results are conditional:
some settlement prices may reflect exchange valuation models,
and the longer study uses other same-date option strikes to
estimate its pricing inputs. The two studies use different
models and must not be pooled as one sample.

The contribution is an empirical comparison of what changes a
WTI average-price option valuation: integrating volatility
uncertainty or changing the volatility input. We also distinguish
uncertainty in model prices from uncertainty in observed pricing
errors. The comparison provides a practical test of whether additional
Bayesian modeling improves conditional pricing. We keep the
limitations of settlement data and model assumptions explicit.
